\documentclass[12pt]{article}
\usepackage{amsmath,amssymb,amsfonts}
\usepackage[margin=1in]{geometry}
\usepackage{bm}

\newcommand{\vect}[1]{\mathbf{#1}}
\newcommand{\grad}{\boldsymbol{\nabla}}

\begin{document}

\section*{Problem 8}

\textbf{Problem:} If the electric potential of a point charge were
\[
\varphi(r) = \frac{q}{r^{1+\varepsilon}}, \qquad \varepsilon \ll 1 \quad \text{(positive or negative)}
\]
instead of $q/r$, many of the results of physics would be altered. Calculate for $r \neq 0$:
\begin{itemize}
\item[(a)] $\vect{E} = -\grad\varphi$
\item[(b)] $\grad\cdot\vect{E} = -\nabla^2\varphi$
\item[(c)] $\grad\times\vect{E}$
\item[(d)] Derive the analog of Gauss's law for a spherical region of radius $R$.
\end{itemize}
Examine the limits as $\varepsilon \to 0$.

\bigskip
\textbf{Solution:}

\subsection*{(a) $\vect{E} = -\grad\varphi$}

Since $\varphi = q\,r^{-(1+\varepsilon)}$ depends only on $r = |\vect{r}|$, the gradient is purely radial:
\[
\vect{E} = -\grad\varphi = -\frac{d\varphi}{dr}\,\hat{\vect{r}} = -q\bigl[-(1+\varepsilon)\,r^{-(2+\varepsilon)}\bigr]\,\hat{\vect{r}}.
\]

\[
\boxed{\vect{E} = \frac{q(1+\varepsilon)}{r^{2+\varepsilon}}\,\hat{\vect{r}}}.
\]

\textit{Limit:} As $\varepsilon \to 0$, $\vect{E} \to q\,\hat{\vect{r}}/r^2$, the standard Coulomb field.

\subsection*{(b) $\grad\cdot\vect{E} = -\nabla^2\varphi$}

For a spherically symmetric function $f(r)$, the Laplacian is:
\[
\nabla^2 f = \frac{1}{r^2}\frac{d}{dr}\!\left(r^2\frac{df}{dr}\right).
\]

Compute:
\[
\frac{d\varphi}{dr} = -(1+\varepsilon)\,q\,r^{-(2+\varepsilon)},
\]
\[
r^2\frac{d\varphi}{dr} = -(1+\varepsilon)\,q\,r^{-\varepsilon},
\]
\[
\frac{d}{dr}\!\left(r^2\frac{d\varphi}{dr}\right) = -(1+\varepsilon)\,q\,(-\varepsilon)\,r^{-\varepsilon-1} = \varepsilon(1+\varepsilon)\,q\,r^{-(1+\varepsilon)}.
\]

Therefore:
\[
\nabla^2\varphi = \frac{\varepsilon(1+\varepsilon)\,q}{r^{3+\varepsilon}},
\]

\[
\boxed{\grad\cdot\vect{E} = -\nabla^2\varphi = -\frac{\varepsilon(1+\varepsilon)\,q}{r^{3+\varepsilon}}}.
\]

\textit{Limit:} As $\varepsilon \to 0$, $\grad\cdot\vect{E} \to 0$ for $r \neq 0$, consistent with the Coulomb result $\nabla^2(1/r) = 0$ for $r \neq 0$.

Note that for $\varepsilon \neq 0$, the divergence is \emph{nonzero} even away from the charge. Gauss's law in its usual differential form fails---the potential $q/r^{1+\varepsilon}$ does not satisfy Laplace's equation.

\subsection*{(c) $\grad\times\vect{E}$}

Since $\vect{E} = -\grad\varphi$ and $\varphi$ is a scalar field, we have identically:
\[
\boxed{\grad\times\vect{E} = -\grad\times(\grad\varphi) = \vect{0}}.
\]

This follows from Eq.~(1.83): $\text{curl}(\text{grad}\,\varphi) = 0$ for any sufficiently smooth scalar field. This result holds for \emph{any} $\varepsilon$ and is not affected by the modification.

\subsection*{(d) Analog of Gauss's law}

Integrate $\grad\cdot\vect{E}$ over a sphere of radius $R$ centered at the origin. Using the divergence theorem:
\[
\oint_S \vect{E}\cdot d\boldsymbol{\sigma} = \int_\tau \grad\cdot\vect{E}\;d\tau.
\]

\textbf{Left side (direct surface integral):}

On the sphere, $\vect{E} = \frac{q(1+\varepsilon)}{R^{2+\varepsilon}}\,\hat{\vect{r}}$ and $d\boldsymbol{\sigma} = R^2\sin\theta\,d\theta\,d\phi\;\hat{\vect{r}}$, so:
\[
\oint_S \vect{E}\cdot d\boldsymbol{\sigma} = \frac{q(1+\varepsilon)}{R^{2+\varepsilon}} \cdot 4\pi R^2 = \frac{4\pi q(1+\varepsilon)}{R^{\varepsilon}}.
\]

\textbf{Right side (volume integral):}

\[
\int_\tau \grad\cdot\vect{E}\;d\tau = -\int_0^R \frac{\varepsilon(1+\varepsilon)\,q}{r^{3+\varepsilon}} \cdot 4\pi r^2\,dr = -4\pi\varepsilon(1+\varepsilon)\,q \int_0^R r^{-(1+\varepsilon)}\,dr.
\]

For $\varepsilon > 0$, this integral diverges at $r = 0$ (reflecting the point-charge singularity). For $\varepsilon < 0$ (with $|\varepsilon| < 1$):
\[
\int_0^R r^{-(1+\varepsilon)}\,dr = \frac{r^{-\varepsilon}}{-\varepsilon}\bigg|_0^R = \frac{R^{-\varepsilon}}{-\varepsilon} - \lim_{r\to 0}\frac{r^{-\varepsilon}}{-\varepsilon}.
\]

For $\varepsilon < 0$, $-\varepsilon > 0$, so $r^{-\varepsilon} \to 0$ as $r \to 0$, giving $\frac{R^{-\varepsilon}}{-\varepsilon}$.

Then:
\[
\int_\tau \grad\cdot\vect{E}\;d\tau = -4\pi\varepsilon(1+\varepsilon)\,q \cdot \frac{R^{-\varepsilon}}{-\varepsilon} = \frac{4\pi(1+\varepsilon)\,q}{R^{\varepsilon}},
\]
which matches the surface integral, confirming consistency.

The \textbf{analog of Gauss's law} for the modified potential is:
\[
\boxed{\oint_S \vect{E}\cdot d\boldsymbol{\sigma} = \frac{4\pi q(1+\varepsilon)}{R^{\varepsilon}}}.
\]

The flux through the sphere now depends on the radius $R$ (unlike the Coulomb case where it equals $4\pi q$ for any $R$). As $\varepsilon \to 0$, $R^{-\varepsilon} \to 1$ and $(1+\varepsilon) \to 1$, recovering $\oint \vect{E}\cdot d\boldsymbol{\sigma} = 4\pi q$. \hfill $\blacksquare$

\end{document}
